\section{Visualisation de textes assistée par LLM}
Story Ribbons extrait avec des LLM des informations narratives de romans et de scripts, puis représente personnages et thèmes à différents niveaux. Les auteurs rapportent des évaluations du pipeline et des études sur 36 œuvres \cite{ref01}. Il montre qu'une extraction par LLM peut alimenter une navigation à plusieurs niveaux. Notre transfert vers l'argumentation reste une hypothèse de conception : la succession d'événements n'est pas une relation de preuve.

Graphologue transforme des réponses de LLM en diagrammes d'entités et de relations, avec des interactions permettant des questions contextualisées \cite{ref02}. Sensecape externalise plusieurs niveaux d'abstraction pour soutenir exploration et organisation des informations \cite{ref03}. Nous en retenons deux exigences : limiter la densité visuelle et garder le contexte local visible quand un élément est sélectionné.

La revue de Brossier et collègues classe 48 articles suivant six dimensions et souligne notamment des enjeux de lecture et d'accessibilité \cite{ref04}. Elle invite à évaluer ensemble l'algorithme et l'interaction : une extraction précise mais difficile à inspecter ne suffit pas, et une interface convaincante peut susciter une confiance que le contenu ne justifie pas.
% TODO: vérifier la publication finale EuroVis/CGF ; la notice consultée indique une soumission.
